% Unit 033 English translation of chapter3.tex lines 163--345.
% Source: Methods of Algebra, Volume 2, commit
% 9a5803ff77dd3257484cb177f851a73770a59dd3 (CC BY 4.0).
% stable-unit-id: o014.aljabr2.chapter3.hom-complex-and-homotopy
% Coverage: Section 3.2 in full; ends before Section 3.3.

% segment-id: o014.li2.chapter3.hom-complex-and-homotopy.h001
\section{\texorpdfstring{$\Hom$}{Hom} Complexes and Homotopy}\label{sec:Hom-cplx}
\hypertarget{o014.aljabr2.chapter3.hom-complex-and-homotopy}{ }

% segment-id: o014.li2.chapter3.hom-complex-and-homotopy.p001
We previously defined $\Hom$ between two complexes. This section explains how
to promote it to a complex and how to promote composition of morphisms to a
kind of multiplication on complexes. We continue to let $\mathcal{A}$ be an
additive category. If $\Bbbk$ is any commutative ring and $\mathcal{A}$ is a
$\Bbbk$-linear category, every statement in this section has an evident
$\Bbbk$-linear version.

% segment-id: o014.li2.chapter3.hom-complex-and-homotopy.p002
Recall that taking $\Hom$ sets gives a functor
$\Hom:\mathcal{A}^{\opp}\times\mathcal{A}\to\cate{Ab}$. The notation $\Hom$
without a superscript refers to the $\Hom$ set in $\mathcal{A}$ or in
$\cate{C}(\mathcal{A})$; when necessary, a subscript distinguishes the
category.

% segment-id: o014.li2.chapter3.hom-complex-and-homotopy.d001
\begin{definition}[$\Hom$ complex]\label{def:Hom-cplx}
	\index[sym1]{Hombullet@$\Hom^{\bullet}$}
	Let $X,Y$ be objects of $\cate{C}(\mathcal{A})$. For every $n\in\Z$, define
% segment-id: o014.li2.chapter3.hom-complex-and-homotopy.q001
	\[ \Hom^n\left(X, Y \right) := \prod_{k \in \Z} \Hom_{\mathcal{A}}\left( X^k, Y^{k+n} \right); \]
	termwise composition of morphisms gives
% segment-id: o014.li2.chapter3.hom-complex-and-homotopy.q002
	\begin{equation}\label{eqn:Hom-cplx-multiplication}\begin{aligned}
		\Hom^n \left(Y, Z \right) \times \Hom^m\left( X, Y \right) & \longrightarrow \Hom^{n+m}\left(X , Z \right) \\
		(g, f) & \longmapsto gf := \left( g^{k+m} f^k \right)_{k \in \Z}.
	\end{aligned}\end{equation}
	Notice that $d_X=(d_X^k)_k\in\Hom^1(X,X)$, and similarly for $d_Y$.
	Accordingly, define
% segment-id: o014.li2.chapter3.hom-complex-and-homotopy.q003
	\begin{align*}
		d^n = d_{\Hom^\bullet(X, Y)}^n: \Hom^n\left(X, Y \right) & \longrightarrow \Hom^{n+1}\left(X, Y \right) \\
		f & \longmapsto d_Y f - (-1)^n f d_X.
	\end{align*}
\end{definition}

% segment-id: o014.li2.chapter3.hom-complex-and-homotopy.p003
The operation \eqref{eqn:Hom-cplx-multiplication} is plainly associative and
distributive over addition. Given morphisms $u:\underline{X}\to X$ and
$v:Y\to\underline{Y}$ in $\cate{C}(\mathcal{A})$, regard them respectively
as elements of $\Hom^0(\underline{X},X)$ and $\Hom^0(Y,\underline{Y})$.
For every $n\in\Z$ there is then a homomorphism
% segment-id: o014.li2.chapter3.hom-complex-and-homotopy.q004
\begin{equation}\label{eqn:Hom-cplx-functoriality}\begin{aligned}
	\Hom^n(u, v): \Hom^n\left( X, Y\right) & \longrightarrow \Hom^n\left(\underline{X}, \underline{Y}\right) \\
	f & \longmapsto vfu.
\end{aligned}\end{equation}

% segment-id: o014.li2.chapter3.hom-complex-and-homotopy.dp001
\begin{definition-proposition}
	\index{Hom complex@$\Hom$ complex}
	The data $\left(\Hom^n(X,Y),d^n\right)_{n\in\Z}$ above form an object of
	$\cate{C}(\cate{Ab})$, called the \emph{$\Hom$ complex}. As $X$ and $Y$
	vary, this construction together with
	\eqref{eqn:Hom-cplx-functoriality} gives an additive functor
% segment-id: o014.li2.chapter3.hom-complex-and-homotopy.q005
	\[ \Hom^\bullet: \cate{C}(\mathcal{A})^{\opp} \times \cate{C}(\mathcal{A}) \to \cate{C}(\cate{Ab}). \]
\end{definition-proposition}
% segment-id: o014.li2.chapter3.hom-complex-and-homotopy.pf001
\begin{proof}
	First verify that $d^{n+1}d^n=0$. Since $d_Y^2=0=d_X^2$, a direct
	calculation gives
% segment-id: o014.li2.chapter3.hom-complex-and-homotopy.q006
	\begin{align*}
		d^{n+1} \left( d^n f \right) & = d_Y \left( d_Y f - (-1)^n f d_X \right) - (-1)^{n+1} \left( d_Y f - (-1)^n f d_X \right) d_X \\
		& = (-1)^{n+1} d_Y f d_X - (-1)^{n+1} d_Y f d_X = 0.
	\end{align*}

	Next consider functoriality. Place the map $\Hom^\bullet(u,v)$ defined in
	\eqref{eqn:Hom-cplx-functoriality} into the diagram%
	\footnote{Translator's note (O014-C033): the source labels both horizontal
	arrows $\Hom^{n+1}(u,v)$. Since the upper arrow connects terms of degree
	$n$, this translation restores its label to $\Hom^n(u,v)$; the lower arrow
	retains the label $\Hom^{n+1}(u,v)$.}
% segment-id: o014.li2.chapter3.hom-complex-and-homotopy.g001
	\[\begin{tikzcd}
		\Hom^n\left( X^\bullet, Y^\bullet \right) \arrow[d, "{d^n}"'] \arrow[r, "{\Hom^n(u, v)}" inner sep=0.7em] & \Hom^n\left( \underline{X}^\bullet, \underline{Y}^\bullet \right) \arrow[d, "{\underline{d}^n}"] \\
		\Hom^{n+1}\left( X^\bullet, Y^\bullet \right) \arrow[r, "{\Hom^{n+1}(u, v)}"' inner sep=0.7em] & \Hom^{n+1}\left( \underline{X}^\bullet, \underline{Y}^\bullet \right) ;
	\end{tikzcd}\]
	Because $d_Xu=ud_{\underline{X}}$ and $vd_Y=d_{\underline{Y}}v$, this
	diagram is plainly commutative.
\end{proof}

% segment-id: o014.li2.chapter3.hom-complex-and-homotopy.p004
The differential $d_{\Hom^\bullet(X,Y)}$ is also compatible with shifts. Its
proof is much easier than its statement.

% segment-id: o014.li2.chapter3.hom-complex-and-homotopy.le001
\begin{lemma}\label{prop:Hom-cplx-d-translate}
	Let $n,m\in\Z$ and $f\in\Hom^n(X,Y)$. Identify $f$ with the element
	$\underline{f}=(f^k)_k$ of $\Hom^{n-m}(X,Y[m])$, or with the element
	$\overline{f}=(f^{k-m})_k$ of $\Hom^{n-m}(X[-m],Y)$. Then
% segment-id: o014.li2.chapter3.hom-complex-and-homotopy.q007
	\begin{align*}
		d^{n-m}_{\Hom^\bullet(X, Y[m])} \underline{f} & \xlongequal{\text{under the identification}} (-1)^m d^n_{\Hom^\bullet(X, Y)} f, \\ d^{n-m}_{\Hom^\bullet(X[-m], Y)} \overline{f} & \xlongequal{\text{under the identification}} d^n_{\Hom^\bullet(X, Y)} f.
	\end{align*}
	Consequently,
	$\Hom^\bullet(X,Y[m])=\Hom^\bullet(X,Y)[m]\simeq
	\Hom^\bullet(X[-m],Y)$.
\end{lemma}
% segment-id: o014.li2.chapter3.hom-complex-and-homotopy.pf002
\begin{proof}
	For example, for the first equality,
% segment-id: o014.li2.chapter3.hom-complex-and-homotopy.q008
	\[ (d^{n-m}_{\Hom^\bullet(X, Y[m])} \underline{f})^k = d_{Y[m]}^{k+n-m} \underline{f}^k - (-1)^{n-m} \underline{f}^{k+1} d_X^k : X^k \to Y^{k+n+1}, \]
	which also equals
	$(-1)^m\left(d_Y^{k+n}f^k-(-1)^nf^{k+1}d_X^k\right)$. The argument for
	$\overline{f}$ is similar: although the case of $\overline{f}$ has no
	factor $(-1)^m$, there is still an isomorphism of $\Hom$ complexes; see
	the footnote to Definition~\ref{def:translation-functor}.
\end{proof}

% segment-id: o014.li2.chapter3.hom-complex-and-homotopy.p005
The multiplication defined in \eqref{eqn:Hom-cplx-multiplication} also obeys
the Leibniz rule.

% segment-id: o014.li2.chapter3.hom-complex-and-homotopy.le002
\begin{lemma}\label{prop:Hom-cplx-Leibniz}
	Let $X,Y,Z$ be complexes and
	$(g,f)\in\Hom^n(Y,Z)\times\Hom^m(X,Y)$. With respect to the
	multiplication in \eqref{eqn:Hom-cplx-multiplication}, the following
	equality holds:
% segment-id: o014.li2.chapter3.hom-complex-and-homotopy.q009
	\[
		\begin{aligned}
			d^{n+m}(gf)
				&= (d^n g)f \\
				&\quad + (-1)^n g(d^m f)
		\end{aligned}
	\]
	in $\Hom^{n+m+1}(X,Z)$.
\end{lemma}
% segment-id: o014.li2.chapter3.hom-complex-and-homotopy.pf003
\begin{proof}
	This is a direct calculation, left to the reader.
\end{proof}

% segment-id: o014.li2.chapter3.hom-complex-and-homotopy.le003
\begin{lemma}\label{prop:Hom-as-cocycle}
	Let $n\in\Z$ and $f\in\Hom^n(X,Y)$. Then $d^nf=0$ if and only if
	$f\in\Hom(X,Y[n])$.
\end{lemma}
% segment-id: o014.li2.chapter3.hom-complex-and-homotopy.pf004
\begin{proof}
	Via Lemma~\ref{prop:Hom-cplx-d-translate}, identify $f$ with an element of
	$\Hom^0(X,Y[n])$, reducing the question to the case $n=0$. Clearly,
	$d^0f=0\iff d_Yf=fd_X$.
\end{proof}

% segment-id: o014.li2.chapter3.hom-complex-and-homotopy.d002
\begin{definition}\label{def:homotopy}
	\index{homotopy}
	\index{null-homotopic}
	Let $n\in\Z$. Regard $\Hom(X,Y[n])$ as a subset of $\Hom^n(X,Y)$.
% segment-id: o014.li2.chapter3.hom-complex-and-homotopy.l001
	\begin{compactitem}
% segment-id: o014.li2.chapter3.hom-complex-and-homotopy.i001
		\item If $f,g\in\Hom(X,Y[n])$ and
		$h\in\Hom^{n-1}(X,Y)$ (equivalently,
		$h\in\Hom^{-1}(X,Y[n])$)%
		\footnote{Translator's note (O014-C034): the source does not state the
		degree of $h$, although $d^{n-1}h$ requires
		$h\in\Hom^{n-1}(X,Y)$. This translation makes its type explicit.}
		satisfies $g-f=d^{n-1}h$, then $h$ is called a \emph{homotopy} from
		$f$ to $g$.
% segment-id: o014.li2.chapter3.hom-complex-and-homotopy.i002
		\item For $f,g$ as above, if there exists a homotopy from $f$ to $g$,
		then $f$ and $g$ are called homotopic. A morphism $f$ homotopic to the
		zero morphism is called \emph{null-homotopic}.
	\end{compactitem}

	As the special case $n=0$, a morphism $f\in\Hom(X,Y)$ is null-homotopic
	if and only if there is a family of morphisms $h^m:X^m\to Y^{m-1}$ such
	that
% segment-id: o014.li2.chapter3.hom-complex-and-homotopy.q010
	\[ \forall m \in \Z, \quad f^m = d_Y^{m-1} h^m + h^{m+1} d_X^m . \]
\end{definition}

% segment-id: o014.li2.chapter3.hom-complex-and-homotopy.p006
Homotopy is an equivalence relation. Lemma~\ref{prop:Hom-as-cocycle} gives the
key equality in $\cate{Ab}$ (or $\Bbbk\dcate{Mod}$):
% segment-id: o014.li2.chapter3.hom-complex-and-homotopy.q011
\[ \Hm^n\left(\Hom^\bullet(X, Y), d^\bullet \right) = \Hom\left(X, Y[n]\right) \big/ \{\text{null-homotopic morphisms} \}. \]

% segment-id: o014.li2.chapter3.hom-complex-and-homotopy.p007
The next result shows that null-homotopic morphisms are not merely closed under
addition but form an ideal.

% segment-id: o014.li2.chapter3.hom-complex-and-homotopy.le004
\begin{lemma}\label{prop:null-homotopic-composition}
	Let $X\xrightarrow{f}Y[m]$ and $Y\xrightarrow{g}Z[n]$ be morphisms of
	complexes. If $f$ or $g$ is null-homotopic, then
	$g[m]\circ f:X\to Z[n+m]$ is null-homotopic.
\end{lemma}
% segment-id: o014.li2.chapter3.hom-complex-and-homotopy.pf005
\begin{proof}
	We suppress the superscripts on $d$. Suppose $f=dh$, where
	$h\in\Hom^{-1}(X,Y[m])$. Interpreting composition inside $\Hom^\bullet$,
	we may abbreviate $g[m]\circ f$ as $gf$. Lemma~\ref{prop:Hom-cplx-Leibniz}
	gives $d(gh)=(dg)h+(-1)^ng(dh)=(-1)^ngf$, so $gf$ is null-homotopic. The
	case $g=dh$ is similar.
\end{proof}

% segment-id: o014.li2.chapter3.hom-complex-and-homotopy.p008
Suppressing the subscript on $d_{\Hom^\bullet}$,
Lemmas~\ref{prop:Hom-cplx-Leibniz} and
\ref{prop:null-homotopic-composition} show that the multiplication on
$\Hom^\bullet$ in \eqref{eqn:Hom-cplx-multiplication} induces
% segment-id: o014.li2.chapter3.hom-complex-and-homotopy.q012
\begin{equation}\label{eqn:homotopy-cat-cmposition}
	\Hm^n\left(\Hom^\bullet(Y, Z) \right) \times \Hm^m\left(\Hom^\bullet(X, Y) \right) \to \Hm^{n+m}\left( \Hom^\bullet(X, Z) \right) ,
\end{equation}
and the binary operation \eqref{eqn:homotopy-cat-cmposition} inherits
associativity and distributivity over addition from
\eqref{eqn:Hom-cplx-multiplication}.

% segment-id: o014.li2.chapter3.hom-complex-and-homotopy.d003
\begin{definition}\label{def:cplx-homotopy-cat}
	\index[sym1]{KA@$\cate{K}(\mathcal{A})$}
	Define the category $\cate{K}(\mathcal{A})$ from
	$\cate{C}(\mathcal{A})$ by setting
% segment-id: o014.li2.chapter3.hom-complex-and-homotopy.q013
	\begin{align*}
		\Obj\left( \cate{K}(\mathcal{A}) \right) & := \Obj\left( \cate{C}(\mathcal{A}) \right), \\
		\Hom_{\cate{K}(\mathcal{A})}\left( X, Y \right) & := \Hm^0\left( \Hom^\bullet\left( X, Y \right), d^\bullet_{\Hom} \right),
	\end{align*}
	with composition determined by \eqref{eqn:homotopy-cat-cmposition} for
	$n=m=0$.
\end{definition}

% segment-id: o014.li2.chapter3.hom-complex-and-homotopy.p009
Observe that $\cate{K}(\mathcal{A})$ is an additive category: addition on
morphism sets is clear, the zero object comes from the zero object of
$\cate{C}(\mathcal{A})$, and the existence of products (or coproducts) follows
from the more general property that there are natural isomorphisms
% segment-id: o014.li2.chapter3.hom-complex-and-homotopy.q014
\begin{align*}
	\Hom_{\cate{K}(\mathcal{A})}\left( X, \prod_{i \in I} Y_i \right) & \simeq \prod_{i \in I} \Hom_{\cate{K}(\mathcal{A})}(X, Y_i) \\
	\text{or}\quad \Hom_{\cate{K}(\mathcal{A})}\left( \coprod_{i \in I} X_i, Y \right) & \simeq \prod_{i \in I} \Hom_{\cate{K}(\mathcal{A})}(X_i, Y),
\end{align*}
where $I$ is any set, provided $\prod_{i\in I}Y_i$ (or $\coprod_iX_i$)
exists at the level of $\cate{C}(\mathcal{A})$. Since taking products and
taking $\Hm^0$ commute in $\cate{C}(\cate{Ab})$, everything reduces to the
evident isomorphisms of $\Hom$ complexes
% segment-id: o014.li2.chapter3.hom-complex-and-homotopy.q015
\begin{align*}
	\Hom^\bullet\left(X, \prod_i Y_i\right) & \simeq \prod_i \Hom^\bullet(X, Y_i) \\
	\text{or}\quad \Hom^\bullet\left(\coprod_i X_i, Y\right) & \simeq \prod_i \Hom^\bullet(X_i, Y).
\end{align*}

% segment-id: o014.li2.chapter3.hom-complex-and-homotopy.p010
The shift functor $[n]$ induces an automorphism of the additive category
$\cate{K}(\mathcal{A})$, still denoted by $[n]$. The additive functor
$\cate{C}(\mathcal{A})\to\cate{K}(\mathcal{A})$ is the identity on objects
and the quotient map on morphism sets. The following universal property is
immediate.

% segment-id: o014.li2.chapter3.hom-complex-and-homotopy.pr001
\begin{proposition}\label{prop:homotopy-category-universal-property}
	If $\mathcal{B}$ is an $\cate{Ab}$-category and the additive functor
	$F:\cate{C}(\mathcal{A})\to\mathcal{B}$ sends every null-homotopic
	morphism to the zero morphism, then $F$ factors uniquely through
	$\cate{C}(\mathcal{A})\to\cate{K}(\mathcal{A})$.
\end{proposition}

% segment-id: o014.li2.chapter3.hom-complex-and-homotopy.p011
The exercises in the next chapter will show that $\cate{K}(\mathcal{A})$ is
generally not an abelian category.

% segment-id: o014.li2.chapter3.hom-complex-and-homotopy.r001
\begin{remark}\index{Hom complex@$\Hom$ complex}\index[sym1]{Homlbullet@$\Hom_\bullet$}
	\label{rem:Hom-chain-cplx}
	For chain complexes (Remark~\ref{rem:cochain-vs-chain}) there is a
	corresponding $\Hom_\bullet$:
% segment-id: o014.li2.chapter3.hom-complex-and-homotopy.q016
	\begin{gather*}
		\Hom_n(X, Y) := \prod_{k \in \Z} \Hom\left(X_k, Y_{k+n} \right), \\
		d_n f = d_Y f - (-1)^n f d_X, \quad f \in \Hom_n(X, Y).
	\end{gather*}
	Indeed, use the convention in Remark~\ref{rem:cochain-vs-chain}, setting
	$X^n=X_{-n}$, $d^n=d_{-n}$, and so on. Then
	$\left(\Hom_n(X,Y),d_n\right)_n$ is precisely the chain complex
	corresponding to $\left(\Hom^n(X,Y),d^n\right)_n$:
% segment-id: o014.li2.chapter3.hom-complex-and-homotopy.q017
	\begin{align*}
		\Hom_n\left(X, Y \right) & = \prod_{k \in \Z} \Hom\left( X_k, Y_{n+k} \right) \xlongequal{h = -k} \prod_{h \in \Z} \Hom\left(X^h, Y^{h-n} \right) \\
		& = \Hom^{-n} \left(X, Y\right),
	\end{align*}
	and the relation between $d_n$ and $d^{-n}$ is similar. Of course, chain
	complexes have a corresponding version of $\cate{K}(\mathcal{A})$ as
	well.
\end{remark}

% segment-id: o014.li2.chapter3.hom-complex-and-homotopy.p012
Given an additive functor $F:\mathcal{A}\to\mathcal{B}$ between additive
categories and any $X,Y\in\Obj(\cate{C}(\mathcal{A}))$, applying $F$
termwise gives a morphism in $\cate{C}(\cate{Ab})$
% segment-id: o014.li2.chapter3.hom-complex-and-homotopy.q018
\begin{align*}
	\Hom^\bullet\left(X, Y\right) & \to \Hom^\bullet\left(\cate{C}F(X), \cate{C}F(Y) \right) \\
	(f^k)_{k \in \Z}\in\Hom^r(X,Y) & \mapsto
	(F(f^k))_{k \in \Z}\in
	\Hom^r\left(\cate{C}F(X),\cate{C}F(Y)\right) \qquad (r\in\Z).
\end{align*}
Taking $\Hm^0$ on both sides gives
\[
	\Hom_{\cate{K}(\mathcal{A})}(X,Y)\longrightarrow
	\Hom_{\cate{K}(\mathcal{B})}
	\bigl(\cate{C}F(X),\cate{C}F(Y)\bigr).
\]
Thus we obtain a functor $\cate{K}F:\cate{K}(\mathcal{A})\to
\cate{K}(\mathcal{B})$ that makes the following diagram commute:
% segment-id: o014.li2.chapter3.hom-complex-and-homotopy.g002
\begin{equation}\label{eqn:KF}\begin{tikzcd}
	\cate{C}(\mathcal{A}) \arrow[r, "{\cate{C}F}"] \arrow[d] & \cate{C}(\mathcal{B}) \arrow[d] \\
	\cate{K}(\mathcal{A}) \arrow[r, "{\cate{K}F}"'] & \cate{K}(\mathcal{B}).
\end{tikzcd}\end{equation}
The top row is the functor in Proposition~\ref{prop:CA-functoriality}. These
constructions are compatible with adjoint pairs.
\index[sym1]{KF@$\cate{K}F$}

% segment-id: o014.li2.chapter3.hom-complex-and-homotopy.pr002
\begin{proposition}\label{prop:KF-adjoint}
	Consider an adjoint pair $(F,G,\varphi)$, where
% segment-id: o014.li2.chapter3.hom-complex-and-homotopy.g003
	$\begin{tikzcd}
		F: \mathcal{A} \arrow[shift left, r] & \mathcal{B} \arrow[shift left, l] : G
	\end{tikzcd}$
	are additive functors between additive categories. Then
% segment-id: o014.li2.chapter3.hom-complex-and-homotopy.g004
	\[\begin{tikzcd}
		\cate{K}F: \cate{K}(\mathcal{A}) \arrow[shift left, r] & \cate{K}(\mathcal{B}) \arrow[shift left, l] : \cate{K}G
	\end{tikzcd}\]
	naturally forms an adjoint pair as well.
\end{proposition}
% segment-id: o014.li2.chapter3.hom-complex-and-homotopy.pf006
\begin{proof}
	The adjunction data $\varphi$ are a family of bijections
	$\varphi_{A,B}:\Hom_{\mathcal{B}}(FA,B)\simeq
	\Hom_{\mathcal{A}}(A,GB)$, functorial in $A$ and $B$.
	Proposition~\ref{prop:automatic-k-morphism} ensures that these bijections
	are linear, and hence they induce a canonical isomorphism of complexes
% segment-id: o014.li2.chapter3.hom-complex-and-homotopy.q019
	\[ \Hom^\bullet\left( \cate{C}F(X), Y \right) \simeq \Hom^\bullet\left( X, \cate{C}G(Y) \right). \]
	In degree $r$, this isomorphism sends
	\[
		(f^k)_{k\in\Z}\in\Hom^r(\cate{C}F(X),Y)
		\longmapsto
		(\varphi_{X^k,Y^{k+r}}(f^k))_{k\in\Z}
		\in\Hom^r(X,\cate{C}G(Y)).
	\]%
	\footnote{Translator's note (O014-C035): the source writes the component
	as $\varphi_{A,B}(f^n)$ without tying $A$ and $B$ to the source and target
	of the relevant graded component. This translation uses $r$ for the
	complex degree and $k$ for the family index, so the component is typed as
	$\varphi_{X^k,Y^{k+r}}(f^k)$.}
	By \eqref{eqn:KF}, taking $\Hm^0$ on both sides makes $\cate{K}F$ left
	adjoint to $\cate{K}G$.
\end{proof}
